\chapter{展望：AI时代的交通工程师应该具备什么能力}

\section{能力重心的迁移}

传统工程师的价值链更强调代码实现、模型调参和报告撰写。AI 时代的工程师价值链将逐渐转向问题定义、Prompting、领域判断、结果评估和方案表达。

领域知识的价值不会因为 AI 降低，反而会提高。AI 最难复制的是对特定场景的深度理解，以及对结果是否可信的判断。

\section{与AI协作的职业伦理}

\begin{enumerate}
  \item \textbf{署名与责任}：即使使用 AI 完成分析，作者仍对最终成果负有完整责任。
  \item \textbf{透明度要求}：在学术论文和工程报告中，应说明 AI 工具的使用方式。
  \item \textbf{不应外包的判断}：涉及安全、政策建议和利益相关方的决策，不能以“AI 这么说”作为免责依据。
  \item \textbf{数据隐私}：将个人位置等敏感数据上传到公共 AI 平台可能带来法律和伦理风险。
\end{enumerate}

\section{持续学习的策略：如何跟上AI工具的快速迭代}

\begin{enumerate}
  \item \textbf{从 Coding 到 Orchestration}：不再自己编写每一行代码，而是组织 AI、数据、模型和验证流程。
  \item \textbf{从 Algorithm 到 Problem Framing}：不只是掌握更多模型，而是更准确地定义问题。
  \item \textbf{从 Result 到 Evidence}：不只展示好看的指标，而是提供可信证据链。
  \item \textbf{从 Individual Skill 到 Human-AI Workflow}：重点不是会不会使用某个 AI，而是能否用 AI 稳定完成复杂项目。
\end{enumerate}
